\chapter{Preface}
\markboth{Preface}{Preface}

\section*{Purpose of This Manual}

This manual tells you how to add QuickEIP to a firmware or desktop project,
how to use its programming interface, and how to connect it to a PLC. It
also tells you how to port QuickEIP to a new platform.

QuickEIP is an EtherNet/IP stack in portable C99. One code base operates as
a \term{scanner}, which sends requests to PLCs and other devices, and as an
\term{adapter}, which a PLC sees as an I/O device.

\section*{Audience}

This manual is for firmware engineers and application programmers. You must
know the C language, the build system of your platform, and the basics of
TCP/IP networks. Knowledge of the Common Industrial Protocol (CIP) helps,
but it is not necessary. The glossary defines the CIP terms that this
manual uses.

\section*{How to Use This Manual}

Table~\ref{tab:roadmap} shows where to start for each task.

\begin{table}[htbp]
\centering
\caption{Where to start}\label{tab:roadmap}
\begin{tabularx}{\linewidth}{@{}P{0.42\linewidth}L@{}}
\toprule
\textbf{Task} & \textbf{Read} \\
\midrule
Learn what QuickEIP does and does not do & Chapter~\ref{ch:intro} \\
Understand the layers and the threading model & Chapter~\ref{ch:arch} \\
Build QuickEIP and set its sizes & Chapter~\ref{ch:build} \\
Make your device an adapter & Chapters~\ref{ch:adapter} and \ref{ch:plc} \\
Read and write PLC tags, or open Class~1 I/O & Chapter~\ref{ch:scanner} \\
Set up NetBurner, ESP-IDF, ClearCore, Linux, or Windows & Chapter~\ref{ch:platforms} \\
Port QuickEIP to a new platform & Chapter~\ref{ch:porting} \\
Find the cause of a fault & Chapter~\ref{ch:trouble} \\
Look up a function, a structure, or a status code & Appendixes~\ref{app:api}--\ref{app:codes} \\
\bottomrule
\end{tabularx}
\end{table}

\section*{Conventions}

\subsection*{Language}

The text of this manual follows ASD-STE100, \textit{Simplified Technical
English}.\footnote{ASD, \textit{ASD-STE100: Simplified Technical English},
issue~8 (Brussels: ASD, 2021).} STE makes the text easy to read and easy to
translate. These are the rules that you will see most:

\begin{itemize}
  \item A procedure step contains one instruction. Each step starts with a
        verb in the imperative, for example ``Set'' or ``Do not.''
  \item A sentence in a procedure has no more than 20 words. A sentence in
        a description has no more than 25 words.
  \item The text uses the active voice and the simple tenses.
  \item A word has one meaning. For example, ``open'' always means ``start
        a session or a connection,'' and ``close'' always means ``stop it.''
  \item A warning or a caution comes before the step that it applies to.
\end{itemize}

Technical names, such as function names, file names, and CIP service names,
are not changed.

\subsection*{Style}

The apparatus of this manual follows \textit{The Chicago Manual of
Style}.\footnote{\textit{The Chicago Manual of Style}, 18th ed. (Chicago:
University of Chicago Press, 2024).} Chapter and section titles use headline
style. Lists of three or more items use the serial comma. Ranges use an en
dash, for example UDP 2223--2231. The manual uses the alternative rule for
numbers in Chicago 9.3: words for zero through nine in running text, and
numerals for 10 and above. Numerals are used for all quantities with a unit,
for data values, for instance numbers, and for code. Notes are footnotes in
the Chicago notes-bibliography style, and the bibliography is at the end of
the manual.

\subsection*{Safety Notices}

\begin{warning}
A warning tells you about a condition that can cause injury or death. It
also tells you about a condition that can make machinery move when you do
not expect it.
\end{warning}

\begin{caution}
A caution tells you about a condition that can cause damage to equipment,
loss of data, or corruption of memory.
\end{caution}

\begin{note}
\notelabel A note gives information that helps you do a task. It does not
tell you about a hazard.
\end{note}

\subsection*{Typography}

\begin{table}[htbp]
\centering
\caption{Typographic conventions}\label{tab:typo}
\begin{tabularx}{\linewidth}{@{}P{0.32\linewidth}L@{}}
\toprule
\textbf{Example} & \textbf{Meaning} \\
\midrule
\code{eip_adapter_start} & A function, a type, a macro, or a value in code. \\
\file{include/eip.h} & A file or a directory, relative to the repository root. \\
\menu{Project Properties} & A menu, a dialog, or a field in a tool. \\
\term{originator} & A term that the glossary defines, at its first use. \\
\hex{01} / \hex{0127} & A CIP general status and its extended status. \\
\OT, \TO & The two directions of a Class 1 connection: originator to target, and target to originator. \\
\bottomrule
\end{tabularx}
\end{table}

All timeouts in the programming interface are in milliseconds, and RPIs are
in microseconds. A function that returns \code{int} returns \code{EIP_OK}
(0) or a negative \code{EIP_ERR_*} code, unless the text gives a different
rule.

\section*{Related Documents}

The comments in the headers in \file{include/} are the primary reference
for the programming interface. Each port with application calls has its own
header: \file{ports/esp-idf/eip_esp_idf.h} and
\file{ports/clearcore/eip_clearcore.h}. The repository also contains
\file{docs/INTEGRATION.md}, \file{docs/PORTING.md}, and
\file{docs/MIGRATION.md}. This manual contains the information in those
files. If this manual and a header comment are different, the header comment
agrees with the code. Use it, and report the difference.

The CIP and EtherNet/IP specifications are published by ODVA. The
bibliography lists them.
